\documentclass[aspectratio=169]{beamer}

% Theme
\usepackage{UFtheme}


% Basic packages
\usepackage{amsmath}
\usepackage{amssymb}
\usepackage{bm}
\usepackage{xcolor}
\usepackage{tikz}
\usetikzlibrary{positioning}

% ------------------------------------------------
% Darker blue for headings / emphasis
% ------------------------------------------------
\definecolor{UFdarkblue}{RGB}{0,0,255}
\definecolor{UFblue}{RGB}{164,3,13}

\setbeamercolor{title}{fg=UFdarkblue}
\setbeamercolor{frametitle}{fg=UFdarkblue}
\setbeamercolor{structure}{fg=UFdarkblue}
\setbeamercolor{alerted text}{fg=UFblue}

% Disable the theme footer globally; re-enable it only on the title slide.
\defbeamertemplate{footline}{titlepage footer}{
    \begin{beamercolorbox}[wd=\paperwidth,ht=2.5ex,dp=1.25ex]{page number in head/foot}
        \usebeamerfont{page number in head/foot}
        \hfill Inversion of Geophysical Data | University of Helsinki | Spring 2026 \hfill \insertframenumber\kern1em
    \end{beamercolorbox}
}
\setbeamertemplate{footline}{}

% Slightly tighter overall line spacing for denser lecture slides
\linespread{0.94}

% Slightly tighten displayed equations and list spacing to reduce slide overflow.
\setlength{\abovedisplayskip}{7pt}
\setlength{\belowdisplayskip}{7pt}
\setlength{\abovedisplayshortskip}{4pt}
\setlength{\belowdisplayshortskip}{4pt}
\addtobeamertemplate{itemize/enumerate body begin}{\setlength{\itemsep}{0.35em}}{}
\addtobeamertemplate{itemize/enumerate subbody begin}{\setlength{\itemsep}{0.2em}}{}

% Avoid PDF viewer artifacts from transparent overlays in presentation mode
\setbeamercovered{invisible}

% ------------------------------------------------
% Minimal lecture-note notation (compile-safe)
% ------------------------------------------------
\newcommand{\mat}[1]{{\mathbf{#1}}}
\renewcommand{\vec}[1]{\mathbf{#1}}

\newcommand{\dd}{ {\vec{d}} }
\newcommand{\dobs}{{\vec{d}^{\text{obs}}}}
\newcommand{\dpre}{{\vec{d}^{\text{pre}}}}
\newcommand{\dmean}{\left\langle \dd \right\rangle}

\newcommand{\mm}{ {\vec{m}} }
\newcommand{\delm}{ {\Delta\mm} }
\newcommand{\mest}{ {\vec{m}^{\text{est}}} }
\newcommand{\mmean}{\left\langle \mm \right\rangle}

% \gg already exists in LaTeX, so renew it
\renewcommand{\gg}[1]{\vec{g}\left(#1\right)}
\newcommand{\ff}[1]{\vec{f}\left(#1\right)}

\newcommand{\GG}{\mat{G}}
\newcommand{\FF}{\mat{F}}
\newcommand{\II}{\mat{I}}

\newcommand{\Cd}{\mat{C}_{\dd}}
\newcommand{\Cm}{\mat{C}_{\mm,\text{prior}}}
\newcommand{\courserepo}{\texttt{https://github.com/pankajkmishra/GEOM2057}}

% ------------------------------------------------
% Document info
% ------------------------------------------------
\title{\LARGE\textbf{Nonlinear Inverse Problems \& Stochastic Inversion}}
\author{Jochen Kamm \& \textcolor{blue}{Pankaj K Mishra}}
\date{\footnotesize\today}

\begin{document}

{
\setbeamertemplate{footline}[titlepage footer]
\begin{frame}[plain]
    \titlepage
\end{frame}
}

% ▰▰▰▰▰▰▰▰▰▰▰▰▰▰▰▰▰▰▰▰▰▰▰▰▰▰▰
\begin{frame}{Learning goals}
\small

By the end of this lecture, the aim is to understand:

\begin{itemize}
    \item[--] what makes a geophysical inverse problem non-linear,
    \item[--] why local methods such as Newton and Gauss--Newton work and where they can fail,
    \item[--] why global and stochastic search methods are useful when the objective surface is complicated,
    \item[--] the basic ideas behind PSO, simulated annealing, and VFSA,
    \item[--] why uncertainty matters and how model ensembles help us think about it.
\end{itemize}

\begin{alertblock}{Intuitive goal}
The main objective is to move from \alert{one best model} to \alert{thinking about search, ambiguity, and uncertainty} in inverse problems.
\end{alertblock}
\end{frame}

% ▰▰▰▰▰▰▰▰▰▰▰▰▰▰▰▰▰▰▰▰▰▰▰▰▰▰▰
\section{RECAP: NL Inversion}


% ------------------------------------------------
% Slide 1: Forward map (linear vs nonlinear)
% ------------------------------------------------
\begin{frame}{What is a non-linear inverse problem? (Forward model)}
\small

\textbf{Forward problem:} predict data from a model $\mm$.

\vspace{0.2cm}

\begin{block}{Linear forward model }
\[
\dpre = \GG\,\mm,
\qquad
\GG\in\mathbb{R}^{N\times M},
\qquad
\dobs = \GG\,\mm + \boldsymbol\varepsilon .
\]
\end{block}

\vspace{0.15cm}

\begin{block}{Non-linear forward model}
\[
\dpre = \gg{\mm},
\qquad
\dobs = \gg{\mm} + \boldsymbol\varepsilon .
\]

\vspace{-0.1cm}
\[
\gg{\alpha\mm_1+\beta\mm_2}\neq \alpha\,\gg{\mm_1}+\beta\,\gg{\mm_2}.
\]
\end{block}

\end{frame}



% ------------------------------------------------
% Slide 2: Inverse problem (objective function)
% ------------------------------------------------
\begin{frame}{What is a non-linear inverse problem? (Inverse task)}
\small

\textbf{Inverse problem:} given observations $\dobs$, estimate $\mm$.

\vspace{0.15cm}

\[
\mest \;=\; \arg\min_{\mm}\ \Phi(\mm)
\]

\vspace{-0.1cm}

\begin{block}{A common objective function}
\[
\Phi(\mm)
=
\underbrace{\left(\dobs-\gg{\mm}\right)^T \Cd^{-1}\left(\dobs-\gg{\mm}\right)}_{\alert{data\ misfit}}
+
\epsilon^2\underbrace{L(\mm)}_{\alert{regularization / prior}}.
\]
\end{block}

{
\vspace{0.1cm}
\begin{itemize}
  \item[--] $\gg{\mm}$ is \alert{non-linear} $\Rightarrow$ the minimizer is usually \alert{not available in closed form}.
  \item[--] We therefore solve the inverse problem \alert{iteratively} (linearize $\gg{\mm}$ locally).
\end{itemize}
}

\end{frame}

% ▰▰▰▰▰▰▰▰▰▰▰▰▰▰▰▰▰▰▰▰▰▰▰▰▰▰▰
\begin{frame}{Why do we need iterative methods?}
\small

Near a current model $\mm_i$, we locally linearize the forward operator:
\[
\gg{\mm_i+\delm}
\approx
\gg{\mm_i} + \GG_i \delm,
\qquad
\GG_i = \left.\frac{\partial \gg{\mm}}{\partial \mm}\right|_{\mm=\mm_i}.
\]

{
\vspace{0.15cm}
\begin{itemize}
    \item[--] A non-linear problem is difficult globally, but often manageable \alert{locally}.
    \item[--] So we solve a sequence of \alert{local approximations}.
    \item[--] At each iteration we:
    \begin{itemize}
        \item[--] linearize around the current model,
        \item[--] compute an update $\delm$,
        \item[--] and move to $\mm_{i+1}=\mm_i+\delm$.
    \end{itemize}
    \item[--] This is why the result can depend strongly on the \alert{starting model}.
\end{itemize}
}
\end{frame}

% ▰▰▰▰▰▰▰▰▰▰▰▰▰▰▰▰▰▰▰▰▰▰▰▰▰▰▰
\begin{frame}{Newton's method: the main idea}
\small

Newton's method approximates the objective function locally by a quadratic:
\[
\Phi(\mm_i+\delm)
\approx
\Phi(\mm_i)
+
\bm{\nabla}\Phi(\mm_i)^T\delm
+
\frac{1}{2}\delm^T \mat{H}(\mm_i)\delm .
\]

The update is then
\[
\delm
=
-\mat{H}(\mm_i)^{-1}\bm{\nabla}\Phi(\mm_i),
\qquad
\mm_{i+1}=\mm_i+\delm.
\]

{
\vspace{0.15cm}
\begin{itemize}
    \item[--] Newton's method uses the local \alert{slope} and \alert{curvature} of the objective.
    \item[--] Intuitively: if the objective looked quadratic near the current point, where would its minimum be?
    \item[--] That gives the next model, and the process is repeated.
\end{itemize}
}
\end{frame}

% ▰▰▰▰▰▰▰▰▰▰▰▰▰▰▰▰▰▰▰▰▰▰▰▰▰▰▰
\begin{frame}{Gradient and Hessian: what do they mean?}
\small

The two key ingredients in Newton's method are
\[
\bm{\nabla}\Phi(\mm)
\qquad \text{and} \qquad
\mat{H}(\mm).
\]

\[
\bm{\nabla}\Phi(\mm)
=
\text{\alert{gradient}},
\qquad
\mat{H}(\mm)
=
\text{\alert{Hessian}}.
\]

{
\vspace{0.15cm}
\begin{itemize}
    \item[--] The \alert{gradient} tells us the local \textbf{direction of change} of the objective.
    \item[--] The \alert{Hessian} tells us the local \textbf{curvature} or shape of the objective surface.
    \item[--] In Newton's method:
    \begin{itemize}
        \item[--] the gradient tells us \emph{which way to move},
        \item[--] the Hessian tells us \emph{how the surface bends}, so it scales the step.
    \end{itemize}
    \item[--] This makes Newton more informative than simple gradient descent.
\end{itemize}
}
\end{frame}

% ▰▰▰▰▰▰▰▰▰▰▰▰▰▰▰▰▰▰▰▰▰▰▰▰▰▰▰
\begin{frame}{Other local optimization methods}
\small

Many local optimizers use the same objective function, but differ in how much derivative information they use.

\begin{block}{Common examples in machine learning}
\begin{itemize}
    \item[--] \textbf{Gradient descent}: uses only $\nabla\Phi(\mm)$ and moves in the steepest downhill direction.
    \item[--] \textbf{Adam}: also uses gradients, but adapts the step size using running averages of past gradients.
    \item[--] \textbf{L-BFGS}: uses gradients and builds a low-memory approximation to curvature, without forming the full Hessian.
\end{itemize}
\end{block}

\begin{alertblock}{Relation to Newton and Gauss--Newton}
These are still \alert{local} optimizers. Newton and Gauss--Newton sit in the same family, but they use curvature information more directly.
\end{alertblock}

\begin{itemize}
    \item[--] First-order methods: gradient descent, Adam.
    \item[--] Quasi-Newton methods: L-BFGS.
    \item[--] Second-order methods: Newton, or approximate second-order methods such as Gauss--Newton.
\end{itemize}
\end{frame}

% ▰▰▰▰▰▰▰▰▰▰▰▰▰▰▰▰▰▰▰▰▰▰▰▰▰▰▰
\begin{frame}{Newton method: high-level algorithm}
\small

\[
\mat{H}(\mm_i)\delm = -\bm{\nabla}\Phi(\mm_i),
\qquad
\mm_{i+1}=\mm_i+\delm.
\]

{
\vspace{0.15cm}
\begin{enumerate}
    \item[--] Choose a starting model $\mm_0$.
    \item[--] Evaluate the forward response $\gg{\mm_i}$.
    \item[--] Compute the gradient $\bm{\nabla}\Phi(\mm_i)$ and Hessian $\mat{H}(\mm_i)$.
    \item[--] Solve for the update $\delm$.
    \item[--] Update the model and repeat until convergence.
\end{enumerate}

\vspace{0.05cm}
\begin{itemize}
    \item[--] Newton can converge very fast near a good solution.
    \item[--] But computing the full Hessian is often costly.
\end{itemize}
}
\end{frame}

% ▰▰▰▰▰▰▰▰▰▰▰▰▰▰▰▰▰▰▰▰▰▰▰▰▰▰▰
\begin{frame}{Gauss--Newton method: A practical version of Newton Method}
\small

Gauss--Newton keeps the Newton idea, but approximates the Hessian by dropping expensive second-derivative terms.

\vspace{0.05cm}

Here $\ff{\mm}$ denotes the model-side penalty or regularization map, and $\FF_i$ is its Jacobian evaluated at $\mm_i$.

\vspace{0.1cm}

A common approximation is
\[
\tilde{\mat{H}}(\mm_i)
\approx
2\left(
\GG_i^T \Cd^{-1}\GG_i
+
\epsilon^2 \FF_i^T\FF_i
\right).
\]

The update becomes
\[
\delm
=
\tilde{\mat{H}}(\mm_i)^{-1}
\left[
\GG_i^T\Cd^{-1}\big(\dd-\gg{\mm_i}\big)
+
\epsilon^2\FF_i^T\big(\ff{\mmean}-\ff{\mm_i}\big)
\right].
\]

{
\vspace{0.15cm}
\begin{itemize}
    \item[--] So \alert{Gauss--Newton is an approximation of Newton's method}.
    \item[--] $\ff{\mm}$ measures how the current model is penalized relative to the prior or regularization target.
    \item[--] $\FF_i = \partial \ff / \partial \mm$ tells us how that penalty changes when the model changes.
    \item[--] It is cheaper and easier to use in practice.
    \item[--] It works well for many \alert{moderately non-linear} problems.
    \item[--] This is why \alert{Gauss--Newton is widely used}.
\end{itemize}
}
\end{frame}

% ▰▰▰▰▰▰▰▰▰▰▰▰▰▰▰▰▰▰▰▰▰▰▰▰▰▰▰
% ▰▰▰▰▰▰▰▰▰▰▰▰▰▰▰▰▰▰▰▰▰▰▰▰▰▰▰
\begin{frame}{Simplified covariance model in the scripts}
\small

The general notation uses covariance matrices $\Cd$ and $\Cm$, while the scripts use a deliberately simpler special case.

\begin{alertblock}{Classroom simplification}
\[
\Cd = \sigma_d^2 \II,
\qquad
\Cm = \II.
\]
So the data misfit in the scripts is
\[
\Phi_d(\mm)
=
\frac{1}{2}\chi_d^2(\mm)
=
\frac{1}{2\sigma_d^2}\left\|\dobs-\gg{\mm}\right\|_2^2.
\]
\end{alertblock}

\begin{itemize}
    \item[--] This keeps the example easy to read and easy to code.
    \item[--] The main concepts are the same: weighted data fit, regularization, and iterative or stochastic search.
    \item[--] Later, the same ideas extend to more realistic covariance models.
\end{itemize}
\end{frame}

% ▰▰▰▰▰▰▰▰▰▰▰▰▰▰▰▰▰▰▰▰▰▰▰▰▰▰▰
\begin{frame}{Code demo: deterministic inversion examples}
\small

The lecture demos for the deterministic part of this lecture are in the repository
\[
\courserepo
\]

\begin{block}{Main files}
\begin{itemize}
    \item[--] \texttt{L6/01-problem-setup-derivatives-etc.jl}: forward model, objective, derivatives.
    \item[--] \texttt{L6/02-data-generation.jl}: synthetic data generation.
    \item[--] \texttt{L6/03-inversion-NewtonMethod.jl}: Newton inversion demo.
    \item[--] \texttt{L6/04-inversion-GaussNewton.jl}: Gauss--Newton inversion demo.
\end{itemize}
\end{block}

\begin{alertblock}{Demo focus}
The deterministic scripts connect the mathematics to the synthetic data, the iteration history, and the final recovered model.
\end{alertblock}
\end{frame}

% ▰▰▰▰▰▰▰▰▰▰▰▰▰▰▰▰▰▰▰▰▰▰▰▰▰▰▰


\begin{frame}{Two practical ways to compute derivatives}
\small

There are several practical ways to obtain
\[
\GG=\frac{\partial \gg{\mm}}{\partial \mm}
\]
are:

{
\vspace{0.15cm}
\textbf{Finite differences}

\begin{itemize}
    \item[--] Easy to understand and implement.
    \item[--] Useful for checking code.
    \item[--] But usually expensive.
\end{itemize}

\vspace{0.1cm}
\textbf{Analytical derivatives}
\begin{itemize}
    \item[--] If the forward map is simple enough, the derivatives can also be calculated analytically.
    \item[--] This is often cleaner and more accurate than finite differences.
\end{itemize}

\vspace{0.1cm}
\textbf{Adjoint method}
\begin{itemize}
    \item[--] Very useful in large PDE-based inverse problems.
    \item[--] More efficient when gradients are needed many times.
    \item[--] We will not go into details here; see the lecture notes.
\end{itemize}
}
\end{frame}

% ▰▰▰▰▰▰▰▰▰▰▰▰▰▰▰▰▰▰▰▰▰▰▰▰▰▰▰
\section{Stochastic Inversion}

\begin{frame}{What changes in the stochastic examples?}
\small

We keep the \alert{same toy inverse problem}, the same forward map, and the same objective structure.

\begin{block}{What stays the same}
\begin{itemize}
    \item[--] the forward model $\gg{\mm}$,
    \item[--] the three-parameter model vector $\mm=[m_1,m_2,m_3]$,
    \item[--] the weighted objective $\Phi(\mm)=\Phi_d(\mm)+\epsilon^2\Phi_m(\mm)$.
\end{itemize}
\end{block}

\begin{alertblock}{What changes}
\begin{itemize}
    \item[--] fewer data are used,
    \item[--] the noise level is higher,
    \item[--] we switch from local derivative-based updates to broader global-search strategies.
\end{itemize}
\end{alertblock}
\end{frame}

\begin{frame}{Why local methods may not be enough}
\small

Local methods such as Newton and Gauss--Newton use \alert{nearby slope and curvature} information.

\vspace{0.15cm}

\begin{itemize}
    \item[--] If the objective surface is smooth and bowl-shaped near the starting model, local methods can be excellent.
    \item[--] But if the objective surface has several valleys, plateaus, or narrow basins, the final answer may depend strongly on the starting model.
    \item[--] In our stochastic version of the problem, fewer data and larger noise make the objective surface less informative and potentially more multi-modal.
\end{itemize}

\begin{alertblock}{Motivation for the next methods}
We now want methods that \alert{explore the search space more broadly}, rather than following only local derivatives from one starting point.
\end{alertblock}
\end{frame}

% ▰▰▰▰▰▰▰▰▰▰▰▰▰▰▰▰▰▰▰▰▰▰▰▰▰▰▰
\begin{frame}{Local minimum versus global minimum}
\small

Let $\mathcal{M}$ denote the admissible model set.

\begin{block}{Local minimum}
\[
\mm^\star \text{ is a local minimum if } \exists\,r>0 \text{ such that }
\Phi(\mm^\star) \leq \Phi(\mm)
\quad \forall\, \mm \in \mathcal{M}\cap B_r(\mm^\star).
\]
\end{block}

\begin{block}{Global minimum}
\[
\mm_{\mathrm{glob}} \text{ is a global minimum if }
\Phi(\mm_{\mathrm{glob}}) \leq \Phi(\mm)
\quad \forall\, \mm \in \mathcal{M}.
\]
\end{block}

\begin{itemize}
    \item[--] A local minimum is best only within a neighborhood.
    \item[--] A global minimum is best over the whole domain.
    \item[--] Global optimisation methods are designed to reduce the chance of getting trapped in a poor local minimum.
\end{itemize}
\end{frame}

% ▰▰▰▰▰▰▰▰▰▰▰▰▰▰▰▰▰▰▰▰▰▰▰▰▰▰▰
\begin{frame}{Grid search: the simplest global search}
\small

\textbf{Intuition:} cover the search domain by a finite grid and test every candidate model.

\vspace{0.15cm}

Suppose each parameter $m_j$ is restricted to an interval
\[
m_j \in [m_j^{\min}, m_j^{\max}], \qquad j=1,\dots,M.
\]

Choose grid values $m_j^{(1)},\dots,m_j^{(K_j)}$ and define
\[
\mathcal{G}
=
\left\{
\mm=(m_1^{(k_1)},\dots,m_M^{(k_M)}) : k_j\in\{1,\dots,K_j\}
\right\}.
\]

Then we compute the objective on every grid point:
\[
\mest = \arg\min_{\mm\in\mathcal{G}} \Phi(\mm),
\qquad
\Phi(\mm)=\Phi_d(\mm)+\epsilon^2\Phi_m(\mm).
\]

\begin{itemize}
    \item[--] This is the most direct way to search globally.
    \item[--] It is extremely useful for teaching because the idea is transparent.
    \item[--] It is also useful for visualizing 1D, 2D, or 3D objective surfaces.
\end{itemize}
\end{frame}

% ▰▰▰▰▰▰▰▰▰▰▰▰▰▰▰▰▰▰▰▰▰▰▰▰▰▰▰
\begin{frame}{Grid search: algorithm and limitation}
\small

\begin{block}{Algorithm}
\begin{enumerate}
    \item[--] Choose lower and upper bounds for every parameter.
    \item[--] Discretize each parameter interval.
    \item[--] Evaluate $\Phi(\mm)$ for all $\mm\in\mathcal{G}$.
    \item[--] Keep the model with the smallest value.
\end{enumerate}
\end{block}

\begin{alertblock}{Main limitation}
The cost grows as $|\mathcal{G}|=\prod_{j=1}^M K_j$, so grid search is excellent for teaching and low-dimensional problems, but not for large parameter spaces.
\end{alertblock}

\begin{itemize}
    \item[--] This rapid growth with dimension is often called the \alert{curse of dimensionality}.
    \item[--] Methods such as PSO and VFSA try to search globally without evaluating every possible grid point.
\end{itemize}
\end{frame}

% ▰▰▰▰▰▰▰▰▰▰▰▰▰▰▰▰▰▰▰▰▰▰▰▰▰▰▰
\begin{frame}{Families of global optimisation methods}
\small

\begin{columns}[T]
\column{0.5\textwidth}
\begin{block}{Population-based methods}
\begin{itemize}
    \item[--] \alert{PSO}
    \item[--] Genetic algorithms
    \item[--] Differential evolution
    \item[--] CMA-ES
\end{itemize}
These methods evolve a \alert{population of candidate models}.
\end{block}

\column{0.5\textwidth}
\begin{block}{Annealing-based methods}
\begin{itemize}
    \item[--] Simulated annealing (SA)
    \item[--] \alert{Very Fast Simulated Annealing (VFSA)}
\end{itemize}
These methods follow a \alert{temperature-controlled stochastic search}.
\end{block}
\end{columns}

\begin{alertblock}{Position in this lecture}
SA accepts occasional uphill moves with a probability controlled by temperature. VFSA is a faster, more efficient variant of SA that still allows large exploratory jumps early in the search.
\end{alertblock}
\end{frame}

% ▰▰▰▰▰▰▰▰▰▰▰▰▰▰▰▰▰▰▰▰▰▰▰▰▰▰▰
\begin{frame}[plain]
\centering
\vfill
{\LARGE\textbf{Particle Swarm Optimisation}}
\vfill
\end{frame}

% ▰▰▰▰▰▰▰▰▰▰▰▰▰▰▰▰▰▰▰▰▰▰▰▰▰▰▰
\section{Particle Swarm Optimisation}

\begin{frame}{PSO: intuition and natural inspiration}
\small

\textbf{Particle Swarm Optimisation (PSO)} is inspired by collective motion in nature, such as
\begin{itemize}
    \item[--] bird flocks,
    \item[--] fish schools,
    \item[--] and other groups that move by combining individual behavior with group information.
\end{itemize}

\vspace{0.1cm}

\begin{alertblock}{Intuition}
Each particle is a candidate model. A particle keeps moving through model space by combining three tendencies:
\begin{itemize}
    \item[--] keep some of its current motion,
    \item[--] return toward the best location it has found itself,
    \item[--] move toward the best location found by the swarm.
\end{itemize}
\end{alertblock}

\begin{itemize}
    \item[--] The method is therefore both \alert{social} and \alert{stochastic}.
    \item[--] It does not require derivatives of the objective function.
\end{itemize}
\end{frame}

% ▰▰▰▰▰▰▰▰▰▰▰▰▰▰▰▰▰▰▰▰▰▰▰▰▰▰▰
\begin{frame}{PSO: mathematical objects}
\small

\begin{columns}[T]
\column{0.54\textwidth}
For particle $p=1,\dots,N_p$ at iteration $k$:
\begin{block}{State variables}
\[
\mm_p^{(k)}: \text{particle position},
\qquad
\vec{v}_p^{(k)}: \text{particle velocity}.
\]
\[
\vec{p}_p^{(k)} = \arg\min_{0\leq s\leq k} \Phi\big(\mm_p^{(s)}\big): \text{personal best},
\]
\[
\mm_g^{(k)} = \arg\min_{p,s} \Phi\big(\mm_p^{(s)}\big): \text{global best}.
\]
\end{block}

\column{0.46\textwidth}
\begin{block}{Meaning of the terms}
\begin{itemize}
    \item[--] $\omega$: inertia weight
    \item[--] $c_1$: pull toward personal best
    \item[--] $c_2$: pull toward global best
    \item[--] $r_1,r_2\sim U(0,1)^M$: random exploration
    \item[--] $\sigma_m$: swarm spread
\end{itemize}
\end{block}
\end{columns}

\begin{alertblock}{Idea}
PSO balances momentum, individual memory, and social information to move a swarm toward promising regions without using derivatives.
\end{alertblock}
\end{frame}

% ▰▰▰▰▰▰▰▰▰▰▰▰▰▰▰▰▰▰▰▰▰▰▰▰▰▰▰
\begin{frame}{PSO: mathematical formulation}
\small

\begin{alertblock}{Main PSO update}
\[
\vec{v}_p^{(k+1)}
=
\omega\,\vec{v}_p^{(k)}
+
 c_1\,r_1^{(k)}\odot\big(\vec{p}_p^{(k)}-\mm_p^{(k)}\big)
+
 c_2\,r_2^{(k)}\odot\big(\mm_g^{(k)}-\mm_p^{(k)}\big),
\]
\[
\mm_p^{(k+1)} = \Pi_{\mathcal{M}}\left(\mm_p^{(k)}+\vec{v}_p^{(k+1)}\right).
\]
\end{alertblock}

Here $\odot$ denotes componentwise multiplication and $\Pi_{\mathcal{M}}$ denotes projection back to the admissible model domain.

\begin{itemize}
    \item[--] The first term keeps momentum from the previous iteration.
    \item[--] The second term pulls the particle toward its own personal best.
    \item[--] The third term pulls the particle toward the swarm best.
\end{itemize}
\end{frame}

% ▰▰▰▰▰▰▰▰▰▰▰▰▰▰▰▰▰▰▰▰▰▰▰▰▰▰▰
\begin{frame}{PSO: algorithmic steps}
\small

\begin{columns}[T]
\column{0.56\textwidth}
\begin{block}{Algorithm}
\begin{enumerate}
    \item[--] Initialize $\mm_p^{(0)}$ and $\vec{v}_p^{(0)}$ inside the bounds.
    \item[--] Evaluate $\Phi(\mm_p^{(0)})$.
    \item[--] Set personal bests $\vec{p}_p^{(0)}$ and swarm best $\mm_g^{(0)}$.
    \item[--] Update particles by the equations above.
    \item[--] Replace personal and global bests whenever an improved objective is found.
    \item[--] Stop after the prescribed number of iterations.
\end{enumerate}
\end{block}

\column{0.44\textwidth}
\begin{alertblock}{Key property}
PSO is a \alert{derivative-free} global optimiser. Exploration comes from the random vectors $r_1,r_2$ and from the fact that many particles search simultaneously.
\end{alertblock}
\end{columns}
\end{frame}

% ▰▰▰▰▰▰▰▰▰▰▰▰▰▰▰▰▰▰▰▰▰▰▰▰▰▰▰
\begin{frame}{What objective do PSO and VFSA minimize?}
\footnotesize

PSO and VFSA use the same objective as the deterministic examples; only the search strategy changes.

If the focus is only on data fit, we usually do not need to write the full $\Phi(\mm)$ explicitly.

\[
\textcolor{gray}{\Phi(\mm)=\Phi_d(\mm)+\epsilon^2\Phi_m(\mm)}
\]

\begin{alertblock}{Weighted data misfit}
\[
\chi_d^2(\mm)
=
\left(\dobs-\gg{\mm}\right)^T\Cd^{-1}\left(\dobs-\gg{\mm}\right).
\]
In the classroom scripts, $\Cd=\sigma_d^2\II$, so
\[
\chi_d^2(\mm)
=
\frac{\left\|\dobs-\gg{\mm}\right\|_2^2}{\sigma_d^2},
\]
and therefore
\[
\Phi_d(\mm)=\frac{1}{2}\chi_d^2(\mm).
\]
\end{alertblock}

\begin{itemize}
    \item[--] So the global-search methods still minimize a noise-weighted fit to the data.
    \item[--] RMS can be reported as a diagnostic, but the search itself is based on the weighted objective.
\end{itemize}
\end{frame}

% ▰▰▰▰▰▰▰▰▰▰▰▰▰▰▰▰▰▰▰▰▰▰▰▰▰▰▰
\begin{frame}{Code demo: PSO example}
\small

The PSO demo is in the repository
\[
\courserepo
\]

\begin{block}{Main files}
\begin{itemize}
    \item[--] \texttt{L6/05-inverse-GaussNewton-difficult.jl}: local Gauss--Newton comparison on the same difficult case.
    \item[--] \texttt{L6/06-data-generation-stochastic.jl}: sparse/noisy settings shared by the comparison scripts.
    \item[--] \texttt{L6/07-inversion-PSO.jl}: difficult non-linear PSO demo.
\end{itemize}
\end{block}

\begin{alertblock}{Demo focus}
The PSO demo illustrates swarm exploration, the difference between current particle quality and global best, and the contraction of the swarm toward good regions.
\end{alertblock}
\end{frame}

% ▰▰▰▰▰▰▰▰▰▰▰▰▰▰▰▰▰▰▰▰▰▰▰▰▰▰▰
\section{VFSA}

\begin{frame}{Simulated annealing: intuition and natural inspiration}
\small

\textbf{Simulated annealing (SA)} is inspired by thermal annealing in materials science.

\begin{itemize}
    \item[--] A material is heated so atoms can move freely.
    \item[--] It is then cooled gradually so the system settles into a low-energy state.
\end{itemize}

\begin{alertblock}{Optimisation analogy}
Replace physical \alert{energy} by the objective function $\Phi(\mm)$ and replace temperature by a control parameter $T$.
\end{alertblock}

\begin{itemize}
    \item[--] High temperature allows broad exploration.
    \item[--] Low temperature encourages settling near good minima.
    \item[--] Unlike deterministic descent, SA can accept some uphill moves.
\end{itemize}
\end{frame}

% ▰▰▰▰▰▰▰▰▰▰▰▰▰▰▰▰▰▰▰▰▰▰▰▰▰▰▰
\begin{frame}[plain]
\centering
\vfill
{\LARGE\textbf{Very Fast Simulated Annealing}}
\vfill
\end{frame}

% ▰▰▰▰▰▰▰▰▰▰▰▰▰▰▰▰▰▰▰▰▰▰▰▰▰▰▰

% ▰▰▰▰▰▰▰▰▰▰▰▰▰▰▰▰▰▰▰▰▰▰▰▰▰▰▰
\begin{frame}{VFSA: proposal mechanism}
\small

\textbf{Very Fast Simulated Annealing (VFSA)} is a special variant of simulated annealing designed to keep the same exploration philosophy while using a faster cooling strategy and efficient long-range perturbations.

\vspace{0.1cm}

VFSA builds a perturbation variable $y_j$ for each parameter direction and uses it to generate the next trial model.

\begin{block}{Heavy-tailed random perturbation}
Choose a random number
\[
u_j \sim U(0,1),
\]
and transform it into a Cauchy-type perturbation
\[
y_j
=
\operatorname{sgn}\!\left(u_j-\frac{1}{2}\right)
T\left[
\left(1+\frac{1}{T}\right)^{|2u_j-1|}-1
\right].
\]
\end{block}

\begin{alertblock}{Intuition}
The variable $y_j$ is heavy-tailed, so most jumps are moderate but occasional large jumps are still possible. This is what helps VFSA escape narrow local minima.
\end{alertblock}

\begin{alertblock}{Role of temperature in generating $y_j$}
When $T$ is high, the perturbation $y_j$ can take larger values, so the method explores distant parts of model space. When $T$ becomes small, the same formula produces shorter perturbations, so the search becomes more local and more careful.
\end{alertblock}
\end{frame}

% ▰▰▰▰▰▰▰▰▰▰▰▰▰▰▰▰▰▰▰▰▰▰▰▰▰▰▰
\begin{frame}{VFSA: from perturbation to trial model}
\small

Once $y_j$ has been generated, the new model is formed componentwise by

\begin{alertblock}{VFSA proposal mechanism}
\[
m_{\mathrm{prop},j}
=
\Pi_{[m_j^{\min},m_j^{\max}]}
\left(
m_{\mathrm{current},j}
+
y_j\big(m_j^{\max}-m_j^{\min}\big)
\right),
\]
\end{alertblock}

Here $\Pi_{[m_j^{\min},m_j^{\max}]}$ enforces the parameter bounds.

\begin{itemize}
    \item[--] $\mm_{\mathrm{current}}$: current model in the random walk.
    \item[--] $u_j$: uniform random number.
    \item[--] $y_j$: heavy-tailed perturbation in the $j$th parameter.
    \item[--] $m_j^{\min},m_j^{\max}$: lower and upper bounds for parameter $j$.
\end{itemize}

\begin{block}{Intuition}
The perturbation $y_j$ first decides \emph{how far to move}. The factor $(m_j^{\max}-m_j^{\min})$ scales that jump to the size of the admissible parameter range, and the projection operator keeps the proposal inside the allowed bounds.
\end{block}
\end{frame}

% ▰▰▰▰▰▰▰▰▰▰▰▰▰▰▰▰▰▰▰▰▰▰▰▰▰▰▰
\begin{frame}{VFSA: acceptance rule}
\small

After proposing a new model, VFSA compares the objective values through
\[
\Delta\Phi = \Phi(\mm_{\mathrm{prop}}) - \Phi(\mm_{\mathrm{current}}).
\]

\begin{alertblock}{Acceptance rule}
\[
P_{\mathrm{acc}} =
\begin{cases}
1, & \Delta\Phi \leq 0,\\
\exp\!\left(-\dfrac{\Delta\Phi}{T}\right), & \Delta\Phi > 0.
\end{cases}
\]
\end{alertblock}

\begin{itemize}
    \item[--] If $\Delta\Phi \leq 0$, the new model improves the objective and is accepted.
    \item[--] If $\Delta\Phi > 0$, the new model is worse, but it may still be accepted.
\end{itemize}

\begin{alertblock}{Role of temperature in the acceptance step}
When $T$ is high, uphill moves are accepted more easily, so the search can escape local minima. When $T$ is low, uphill moves are rarely accepted, so the search focuses on refinement near good solutions.
\end{alertblock}
\end{frame}

% ▰▰▰▰▰▰▰▰▰▰▰▰▰▰▰▰▰▰▰▰▰▰▰▰▰▰▰
\begin{frame}{VFSA: cooling schedule}
\small

The temperature is reduced gradually according to a cooling law.

\begin{alertblock}{Cooling schedule}
\[
T_k = \max\!\left(\frac{T_0}{1+0.15(k-1)},10^{-3}\right).
\]
\end{alertblock}

\begin{itemize}
    \item[--] At the beginning, $T_k$ is large, so proposals are bold and uphill moves are still possible.
    \item[--] As $k$ increases, $T_k$ decreases, so both the perturbations and the acceptance of worse moves become smaller.
    \item[--] The search therefore evolves from broad exploration to local refinement.
\end{itemize}

\begin{block}{Intuition}
Temperature is the mechanism that controls the transition from \alert{searching widely} to \alert{settling carefully}.
\end{block}
\end{frame}

% ▰▰▰▰▰▰▰▰▰▰▰▰▰▰▰▰▰▰▰▰▰▰▰▰▰▰▰
\begin{frame}{VFSA: algorithmic steps}
\small

\begin{columns}[T]
\column{0.56\textwidth}
\begin{block}{Algorithm}
\begin{enumerate}
    \item[--] Choose an initial model and initial temperature $T_0$.
    \item[--] For each temperature level, generate several trial moves.
    \item[--] Compute $\Delta\Phi$ and accept with probability $P_{\mathrm{acc}}$.
    \item[--] Update the best-so-far model whenever improvement occurs.
    \item[--] Lower the temperature and continue.
\end{enumerate}
\end{block}

\column{0.44\textwidth}
\begin{alertblock}{Key property}
VFSA is a \alert{single-chain derivative-free} global search method. It explores by random proposals and temperature-controlled acceptance rather than by gradients or swarm motion.
\end{alertblock}
\end{columns}
\end{frame}

% ▰▰▰▰▰▰▰▰▰▰▰▰▰▰▰▰▰▰▰▰▰▰▰▰▰▰▰
\begin{frame}{Code demo: VFSA example}
\small

The VFSA demo is in the repository
\[
\courserepo
\]

\begin{block}{Main files}
\begin{itemize}
    \item[--] \texttt{L6/05-inverse-GaussNewton-difficult.jl}: local Gauss--Newton comparison on the same difficult case.
    \item[--] \texttt{L6/06-data-generation-stochastic.jl}: sparse/noisy settings shared by the comparison scripts.
    \item[--] \texttt{L6/08-inversion-VFSA.jl}: difficult non-linear VFSA demo.
\end{itemize}
\end{block}

\begin{alertblock}{Demo focus}
The VFSA demo illustrates temperature, random proposals, occasional uphill acceptance, and the difference between broad early exploration and later refinement.
\end{alertblock}
\end{frame}

% ▰▰▰▰▰▰▰▰▰▰▰▰▰▰▰▰▰▰▰▰▰▰▰▰▰▰▰
\section{Uncertainty}

\begin{frame}{Why uncertainty matters in geophysical inversion}
\small

In geophysical inversion, a single best-fit model is rarely the whole story.

\begin{alertblock}{Main idea}
Different models can explain the observed data almost equally well, especially when the data are noisy, sparse, or insensitive to some parameters.
\end{alertblock}

\begin{itemize}
    \item[--] We want not only an estimate $\mest$, but also an idea of \alert{how certain} or \alert{how uncertain} that estimate is.
    \item[--] Uncertainty tells us which features of the model are strongly constrained and which are ambiguous.
    \item[--] This matters for interpretation, decision making, and planning further data acquisition.
\end{itemize}
\end{frame}

% ▰▰▰▰▰▰▰▰▰▰▰▰▰▰▰▰▰▰▰▰▰▰▰▰▰▰▰
\begin{frame}{What creates uncertainty?}
\small

\begin{columns}[T]
\column{0.5\textwidth}
\begin{block}{Data-side uncertainty}
\begin{itemize}
    \item[--] measurement noise,
    \item[--] incomplete spatial coverage,
    \item[--] limited resolution,
    \item[--] inconsistent or biased observations.
\end{itemize}
\end{block}

\column{0.5\textwidth}
\begin{block}{Model-side uncertainty}
\begin{itemize}
    \item[--] non-uniqueness,
    \item[--] regularization choices,
    \item[--] parameterization choices,
    \item[--] imperfect forward physics.
\end{itemize}
\end{block}
\end{columns}

\begin{alertblock}{Consequence}
Even if we find a model with small $\Phi(\mm)$, there may still be a broad family of acceptable models around it.
\end{alertblock}
\end{frame}

% ▰▰▰▰▰▰▰▰▰▰▰▰▰▰▰▰▰▰▰▰▰▰▰▰▰▰▰
\begin{frame}{How do we describe uncertainty in this lecture?}
\small

In this lecture, we use a practical viewpoint: uncertainty means that there is not just one acceptable model, but a \alert{family of acceptable models}.

\begin{itemize}
    \item[--] Some model parameters may be tightly constrained, while others may vary a lot with little change in data fit.
    \item[--] We can therefore talk about model spread, acceptable ranges, and families of plausible solutions.
    \item[--] This gives a useful uncertainty picture even before we introduce formal sampling theory.
\end{itemize}

\begin{block}{Practical summary tools}
We summarize uncertainty through model ensembles, spreads, variability of parameters, and sensitivity of the solution to starting models and noise.
\end{block}
\end{frame}

% ▰▰▰▰▰▰▰▰▰▰▰▰▰▰▰▰▰▰▰▰▰▰▰▰▰▰▰
\begin{frame}{Why do we rarely do this fully in 3D?}
\small

\begin{alertblock}{Core difficulty}
In 3D inversion, the number of unknowns is often extremely large, so uncertainty analysis becomes much more expensive than finding one best-fit model.
\end{alertblock}

\begin{itemize}
    \item[--] A 3D model may contain thousands to millions of parameters.
    \item[--] Each forward simulation can already be expensive.
    \item[--] Stochastic uncertainty analysis requires many forward evaluations, not just one inversion run.
    \item[--] Storage and analysis of many 3D models can also become a serious computational burden.
\end{itemize}
\end{frame}

% ▰▰▰▰▰▰▰▰▰▰▰▰▰▰▰▰▰▰▰▰▰▰▰▰▰▰▰
\begin{frame}{What is a model ensemble?}
\small

A \alert{model ensemble} is a collection of acceptable models rather than a single estimate:
\begin{alertblock}{Ensemble}
\[
\mathcal{E} = \left\{ \mm^{(1)},\mm^{(2)},\dots,\mm^{(N)} \right\}.
\]
\end{alertblock}

\begin{itemize}
    \item[--] Each member of the ensemble is a plausible model according to the data and the inversion assumptions.
    \item[--] The spread of the ensemble gives a practical picture of uncertainty.
    \item[--] Ensemble statistics can summarize central tendency and variability.
\end{itemize}

\begin{block}{Simple ensemble summaries}
\[
\bar{\mm} = \frac{1}{N}\sum_{i=1}^{N} \mm^{(i)},
\qquad
\sigma_{m,j}^2 = \frac{1}{N}\sum_{i=1}^{N}\left(m_j^{(i)}-\bar{m}_j\right)^2.
\]
\end{block}
\end{frame}

% ▰▰▰▰▰▰▰▰▰▰▰▰▰▰▰▰▰▰▰▰▰▰▰▰▰▰▰
\begin{frame}{How can VFSA produce a model ensemble?}
\small

One practical approach is to run \alert{multiple independent VFSA chains}.

\begin{block}{Basic idea}
\begin{enumerate}
    \item[--] Start VFSA from many different initial models.
    \item[--] Run each chain with independent random perturbations.
    \item[--] Keep either the final accepted models or a set of low-objective models from each chain.
    \item[--] Merge these models into one ensemble.
\end{enumerate}
\end{block}

\begin{alertblock}{Ensemble from multiple chains}
\[
\mathcal{E} = \bigcup_{r=1}^{R} \left\{ \mm_{r}^{(k)} : k \in \mathcal{K}_{r} \right\},
\]
where $r$ indexes the chain and $\mathcal{K}_{r}$ denotes the saved states from chain $r$.
\end{alertblock}
\end{frame}

% ▰▰▰▰▰▰▰▰▰▰▰▰▰▰▰▰▰▰▰▰▰▰▰▰▰▰▰
\begin{frame}{What do we learn from such an ensemble?}
\small

\begin{itemize}
    \item[--] If all chains converge to a narrow cluster of similar models, uncertainty is relatively small.
    \item[--] If different chains produce very different but still acceptable models, uncertainty is large and non-uniqueness is important.
    \item[--] We can inspect marginal spreads, correlations between parameters, and robust versus unstable model features.
\end{itemize}

\begin{alertblock}{Important caution}
An ensemble built from VFSA runs is a practical uncertainty tool, but its quality still depends on proposal design, cooling schedule, run length, and selection criteria.
\end{alertblock}
\end{frame}

% ▰▰▰▰▰▰▰▰▰▰▰▰▰▰▰▰▰▰▰▰▰▰▰▰▰▰▰
\begin{frame}{Code demo: ensemble idea from multiple VFSA chains}
\small

The uncertainty extension can also be built from the same repository
\[
\courserepo
\]

\begin{block}{Demo focus}
\begin{itemize}
    \item[--] Several VFSA runs can start from different initial models.
    \item[--] Final or low-objective models from each run can be collected.
    \item[--] The spread of accepted models across runs can then be compared.
\end{itemize}
\end{block}

\begin{alertblock}{Main message}
This turns the optimisation demo into a first uncertainty demo: not one answer, but a family of plausible models with visible spread.
\end{alertblock}
\end{frame}

% ▰▰▰▰▰▰▰▰▰▰▰▰▰▰▰▰▰▰▰▰▰▰▰▰▰▰▰
\begin{frame}{Bridge to the next lecture: MCMC}
\small

Today we used global search and ensembles to build intuition about uncertainty. In the next lecture, we will make this more formal using \alert{MCMC} methods.

\vspace{0.1cm}

\begin{itemize}
    \item[--] VFSA gives us a practical way to explore multiple acceptable models.
    \item[--] MCMC will take the next step and treat uncertainty through a more systematic probabilistic sampling framework.
    \item[--] So the current lecture builds the intuition, and the next lecture provides the formal machinery.
\end{itemize}
\end{frame}

% ▰▰▰▰▰▰▰▰▰▰▰▰▰▰▰▰▰▰▰▰▰▰▰▰▰▰▰
\section{Conclusions}

\begin{frame}{So what did we learn?}
\small

{\setlength{\itemsep}{0.2em}
\begin{itemize}
    \item[--] A \alert{non-linear inverse problem} has a forward operator $\dd=\gg{\mm}$ that is not globally linear.
    \item[--] Because of that, we usually solve it \alert{iteratively}.
    \item[--] \alert{Newton's method} uses gradient and Hessian to minimize a local quadratic approximation.
    \item[--] \alert{Gauss--Newton} makes Newton practical by approximating the Hessian.
    \item[--] Deterministic inversion is usually fast and informative, but it can depend strongly on the starting model and may miss distant good solutions.

    \item[--] Stochastic methods such as PSO and VFSA explore more broadly and can reveal multiple acceptable models.
    \item[--] They are powerful for complex objective landscapes, but they are also more computationally expensive.
    \item[--] So deterministic and stochastic inversion each have their own \alert{advantages and limitations}.
\end{itemize}

\vspace{0.12cm}

After this lecture, you should be able to explain the difference between local and global search, describe the ideas behind PSO and VFSA, and interpret uncertainty in terms of families of acceptable models rather than one single answer.
}
\end{frame}

\end{document}